\documentclass{article}
\usepackage[utf8]{inputenc}
\usepackage{amsmath}
\usepackage{amssymb}
\usepackage{stmaryrd}
\usepackage[left=2cm, right=2cm]{geometry}

\title{CS 6371 - Assignment 6}
\author{Christian Parker \\ CXP220034@utdallas.edu}

\begin{document}

\maketitle
\begin{enumerate}
    \item % Question 1
    \begin{enumerate}
        \item % Question 1a
        \begin{equation}
            \begin{aligned}
                e &= (or\ (not\ true)\ false) \\
                &= ((\lambda a.\lambda b.(a\ ?\ true : b))\ (not\ true)\ false) \\
                &= ((\lambda a.\lambda b.((a)(true)(b)))\ (not\ true)\ false) \\
                &= ((\lambda a.\lambda b.((a)(true)(b)))\ ((\lambda b.(b\ ?\ false : true))\ true)\ false) \\
                &= ((\lambda a.\lambda b.((a)(true)(b)))\ ((\lambda b.((b)(false)(true)))\ true)\ false) \\
                &= ((\lambda a.\lambda b.((a)(\lambda x.\lambda y.x)(b)))\ ((\lambda b.((b)(false)(\lambda x.\lambda y.x)))\ (\lambda x.\lambda y.x))\ false) \\
                &= ((\lambda a.\lambda b.((a)(\lambda x.\lambda y.x)(b)))\ ((\lambda b.((b)(\lambda x.\lambda y.y)(\lambda x.\lambda y.x)))\ (\lambda x.\lambda y.x))\ (\lambda x.\lambda y.y)) \\
            \end{aligned}
        \end{equation}

        \item % Question 1b
        \begin{equation}
            \begin{aligned}
                &((\lambda a.\lambda b.((a)(\lambda x.\lambda y.x)(b)))
                    \ ((\lambda b.((b)(\lambda x.\lambda y.y)(\lambda x.\lambda y.x)))
                        (\lambda x.\lambda y.x))
                    \ (\lambda x.\lambda y.y)) \\
                \rightarrow_{1} &((\lambda b.(
                        ((\lambda b.((b)(\lambda x.\lambda y.y)(\lambda x.\lambda y.x)))
                            (\lambda x.\lambda y.x))
                        (\lambda x.\lambda y.x)(b)))
                    \ (\lambda x.\lambda y.y)) \\
                \rightarrow_{1} &(((\lambda b.((b)(\lambda x.\lambda y.y)(\lambda x.\lambda y.x)))
                            (\lambda x.\lambda y.x))
                        \ (\lambda x.\lambda y.x)
                        \ (\lambda x.\lambda y.y)) \\
                \rightarrow_{1} &(((
                            (\lambda x.\lambda y.x)
                            (\lambda x.\lambda y.y)
                            (\lambda x.\lambda y.x))
                        \ (\lambda x.\lambda y.x))
                        \ (\lambda x.\lambda y.y)) \\
                \rightarrow_{1} &(((
                            (\lambda y.(\lambda x.\lambda y.y))
                            (\lambda x.\lambda y.x))
                        \ (\lambda x.\lambda y.x))
                        \ (\lambda x.\lambda y.y)) \\
                \rightarrow_{1} &((
                            (\lambda x.\lambda y.y)
                        \ (\lambda x.\lambda y.x))
                        \ (\lambda x.\lambda y.y)) \\
                \rightarrow_{1} &((\lambda y.y)
                        \ (\lambda x.\lambda y.y)) \\
                \rightarrow_{1} &(\lambda x.\lambda y.y) \\
            \end{aligned}
        \end{equation}
        
    \end{enumerate}

    \item % Question 2
    \begin{enumerate}
        \item % Question 2a
        \begin{equation}
            \begin{aligned}
                isempty = \lambda \ell.\ not\ (\pi_1\ \ell) \\
            \end{aligned}
        \end{equation}

        \item % Question 2b
        \begin{equation}
            \begin{aligned}
                head = \lambda \ell.\ \pi_1\ (\pi_2\ \ell) \\
                tail = \lambda \ell.\ \pi_2\ (\pi_2\ \ell) \\
            \end{aligned}
        \end{equation}

        \item % Question 2c
        \begin{equation}
            \begin{aligned}
                cons = \lambda e. \lambda \ell.(pair\ true\ (pair\ e\ \ell)) \\
            \end{aligned}
        \end{equation}

        \item % Question 2d
        \begin{equation}
            \begin{aligned}
                length = Y(\lambda f. \lambda \ell.(isempty\ \ell\ ?\ 0_\mathbb{N} : succ_\mathbb{N} (f\ (tail\ \ell)))) \\
            \end{aligned}
        \end{equation}

        \item % Question 2e
        \begin{equation}
            \begin{aligned}
                rev &= Y(\lambda f.(\lambda a. \lambda \ell.(isempty\ \ell\ ?\ a : f\ ((head\ \ell) :: a)\ (tail\ \ell))))\ ([]) \\
                append\_helper &= Y(\lambda f.(\lambda \ell_2. \lambda \ell_1.(isempty\ \ell_1\ ?\ \ell_2 : f\ ((head\ \ell_1) :: \ell_2)\ (tail\ \ell_1)))) \\
                append &= \lambda \ell_1. \lambda \ell_2. (append\_helper\ \ell_2\ (rev\ \ell_1)) \\
            \end{aligned}
        \end{equation}

        \item % Question 2f
        \begin{equation}
            \begin{aligned}
                map\_helper &= Y(\lambda m. \lambda a. \lambda f. \lambda \ell.(isempty\ \ell\ ?\ a : m\ ((f\ (head\ \ell)) :: a)\ f\ (tail\ \ell)))\ ([]) \\
                map &= \lambda f. \lambda \ell.\ rev\ ( map\_helper\ f\ \ell ) \\
            \end{aligned}
        \end{equation}

        \item % Question 2g
        \begin{equation}
            \begin{aligned}
                fold\_left &= Y(\lambda g. \lambda f. \lambda a. \lambda \ell.(isempty\ \ell\ ?\ a : g\ (f)\ (f\ a\ (head\ \ell))\ (tail\ \ell))) \\
            \end{aligned}
        \end{equation}
    \end{enumerate}

    \item % Question 3
    \textbf{Theorem.} If \(e \Downarrow e''\) and \(e \rightarrow_1 e'\) then \(e' \Downarrow e''\).
    \\~\\
    \textbf{Proof.}
    There exists a derivation \(D_L\) of \(e \Downarrow e''\) and a derivation \(D_S\) of
    \(e \rightarrow_1 e'\).
    \\~\\
    We will prove \(e' \Downarrow e''\) by structural induction on \(D_S\).
    \\~\\
    \textbf{Inductive Hypothesis.} If \(e_0 \Downarrow e_0''\) has a derivation smaller than \(D_L\)
    and \(e_0 \rightarrow_1 e_0'\) has a derivation smaller than \(D_S\),
    then \(e_0' \Downarrow e_0''\).
    \\~\\
    \textbf{Inductive Cases.}

    \begin{enumerate}
        \item[]\textbf{Case 1.}
        Suppose \(D_S\) ends in S1.
        Then \(e = (\lambda v.e_1) e_2\) and \(e' = e_1[e_2/v]\).

        We must show \(e_1[e_2/v] \Downarrow e''\).
        \\~\\
        Since \(e = (\lambda v.e_1) e_2\), \(D_L\) must have the following form:

        \begin{equation}
            \begin{aligned}
                D_L = \cfrac{\lambda v.e_1 \Downarrow \lambda v.e_1 \qquad e_1[e_2/v] \Downarrow e''}
                {(\lambda v.e_1) e_2 \Downarrow e''}
            \end{aligned}
        \end{equation}
        \(e_1[e_2/v] \Downarrow e''\) is shown by \(D_L\).
        \\~\\
        \item[]\textbf{Case 2.}
        Suppose \(D_S\) ends in S2.
        Then \(D_S\) has the form:
        \begin{equation}
            D_S = \cfrac{\cfrac{D_0}{e_1 \rightarrow_1 e_1'}}{e_1 e_2 \rightarrow_1 e_1' e_2} \\
        \end{equation}

        We must show \(e_1' e_2 \Downarrow e''\).
        Since \(e = e_1 e_2\), \(D_L\) must have the following form:

        \begin{equation}
            D_L = \cfrac{\cfrac{D_1}
                        {e_1 \Downarrow \lambda v.e} \qquad \cfrac{D_2}
                        {e [e_2 / v] \Downarrow e''}}
                    {e_1 e_2 \Downarrow e''} \\
        \end{equation}

        Observe that \(e_1' e_2 \Downarrow e''\) has derivation:
        \\~\\
        \begin{equation}
            \cfrac{\cfrac{D_1'}{e_1' \Downarrow \lambda v.e} \qquad \cfrac{D_2}{e[e_2/v] \Downarrow e''}}
                {e_1' e_2 \Downarrow e''}
        \end{equation}

        We apply the inductive hypothesis with \(D_0 < D_S\), \(D_1 < D_L\), \(e_0 = e_1\),
        \(e_0' = e_1'\), and \(e_0'' = \lambda v.e\).

        Therefore, \(e_1' \Downarrow \lambda v.e\), which is \(D_1'\).
        We conclude \(e_1' e_2 \Downarrow e''\) from its derivation. \(\square\)
    \end{enumerate}
\end{enumerate}

\end{document}